\documentclass[11pt]{scrartcl}
\usepackage[mdthm,colorsec,noauthor,noasy]{evan}

\DeclareMathOperator{\nul}{null}
\DeclareMathOperator{\range}{range}
\DeclareMathOperator{\spn}{span}
\newcommand{\LL}{\mathcal L}
\newcommand{\PP}{\mathcal P}
\newcommand{\MM}{\mathcal M}
\newcommand{\Fp}{\mathbb F_p}
\newcommand{\traps}{\par\medskip\noindent\textbf{\color{RawSienna}Traps.}\ }
\setlength{\emergencystretch}{3em}

\title{Math 110 Quiz Prep}
\subtitle{Notes Ch.~1--9, HW1--4}
\author{}
\date{}

\ihead{}
\begin{document}
\maketitle

\section{Review}

\subsection{Sets and functions}
Let $f\colon X\to Y$ be a function.
For a subset $U$ of $X$, the \vocab{image} $f(U)$ is the set of all outputs $f(x)$ with $x$ in $U$.
For a subset $V$ of $Y$, the \vocab{preimage} $f^{-1}(V)$ is the set of all inputs $x$ whose output $f(x)$ lands in $V$.
The preimage makes sense for every function. Writing $f^{-1}(V)$ does not mean $f$ has an inverse.

\begin{proposition*}[HW1 P1]
  Going to the preimage and back never gives more than you started with: $f(f^{-1}(V))\subseteq V$.
  Going to the image and back never gives less: $U\subseteq f^{-1}(f(U))$.
  \begin{itemize}
    \ii The first is always an equality exactly when $f$ is surjective.
    \ii The second is always an equality exactly when $f$ is injective.
  \end{itemize}
\end{proposition*}

\traps Do not assume you get back what you started with. Take $f(x)=x^2$ on $\RR$.
With $U=[0,1]$: $f(U)=[0,1]$, and $f^{-1}([0,1])=[-1,1]$, which is bigger than $U$.
With $V=[-1,1]$: $f^{-1}(V)=[-1,1]$, and $f([-1,1])=[0,1]$, which is smaller than $V$.

\subsection{Fields}
A \vocab{field} is a set of numbers where you can add, subtract, multiply, and divide by anything except $0$, with the usual rules of algebra.
Examples: $\QQ$, $\RR$, $\CC$, and $\Fp=\{0,1,\dots,p-1\}$ with arithmetic mod $p$, when $p$ is prime.

\begin{proposition*}[HW1 P3]
  In a field, if $ab=0$ then $a=0$ or $b=0$.
\end{proposition*}
\begin{proof}
  Suppose $a\ne0$. Then $a$ has an inverse. Multiply both sides of $ab=0$ by $a^{-1}$ to get $b=0$.
\end{proof}
For example, arithmetic mod $6$ is not a field, because $2\cdot3=6=0$ but neither $2$ nor $3$ is $0$.

In $\FF_2$ we have $1+1=0$. So $(x+1)^2=x^2+2x+1=x^2+1$, and squaring twice more gives $(x+1)^8=x^8+1$ (HW1 P7).

\traps Any step where you divide by $2$ breaks in $\FF_2$, because $2=0$ there.
Over a finite field, two different polynomials can give the same function. For example, $x^2$ and $x$ give the same output at both points of $\FF_2$.

\subsection{Vector spaces and subspaces}
Examples of vector spaces: $\FF^n$; $\FF^\infty$ (infinite sequences); functions from a set $S$ to $\FF$; $\PP_n(\FF)$ (polynomials of degree at most $n$); $\FF^{m,n}$ ($m\times n$ matrices); $\LL(V,W)$ (linear maps from $V$ to $W$).
Basic facts that follow from the axioms: $0v=0$, $a0=0$, $(-1)v=-v$, and if $av=0$ then $a=0$ or $v=0$ (HW1 P4).

\begin{definition*}
  A \vocab{subspace} of $V$ is a subset $U$ that contains $0$, and is closed under addition and scalar multiplication.
  A quick way to check both closure conditions at once: $\lambda u+w$ is in $U$ whenever $u,w$ are in $U$.

  The \vocab{sum} $U+W$ is the set of all vectors $u+w$ with $u$ in $U$ and $w$ in $W$. It is the smallest subspace containing both.

  The sum is \vocab{direct}, written $U\oplus W$, if every vector in $U+W$ can be written as $u+w$ in only one way.
\end{definition*}

\begin{proposition*}[LADR 1.45]
  For two subspaces: $U+W$ is direct exactly when the only vector in both $U$ and $W$ is $0$.

  For any number of subspaces: $U_1+\dots+U_m$ is direct exactly when the only way to write $0=u_1+\dots+u_m$ is with every $u_i=0$.
\end{proposition*}

\traps With three or more subspaces, checking that each pair meets only at $0$ is not enough. Three different lines through the origin in $\RR^2$ pass that check, but their sum is not direct.

The union $U\cup W$ is almost never a subspace. It is one only when one of $U$, $W$ contains the other.

Over $\RR$, a set that is nonempty and closed under addition need not be a subspace (the integer points $\ZZ^n$). Over $\Fp$ it always is, because every scalar is $1+1+\dots+1$.

\subsection{Span, independence, bases}
\begin{definition*}
  The \vocab{span} of $v_1,\dots,v_m$ is the set of all linear combinations $a_1v_1+\dots+a_mv_m$.

  The list is \vocab{linearly independent} if the only combination equal to $0$ is the one with all coefficients $0$.

  A \vocab{basis} is a list that is independent and spans $V$. Equivalently, every vector in $V$ can be written as a combination of the list in exactly one way.
\end{definition*}

\begin{proposition*}[Linear dependence lemma, LADR 2.19]
  If a list is dependent, then one of its vectors is a combination of the earlier ones, and you can delete it without changing the span.
\end{proposition*}

\begin{theorem*}[LADR 2.22]
  In a finite-dimensional space, an independent list is never longer than a spanning list.
\end{theorem*}

\begin{proposition*}[LADR 2.30, 2.32]
  You can delete vectors from any spanning list to get a basis.
  In a finite-dimensional space, you can add vectors to any independent list to get a basis.
\end{proposition*}

\traps A span only contains \emph{finite} combinations. So $e_1,e_2,e_3,\dots$ do not span $\FF^\infty$: the sequence $(1,1,1,\dots)$ would need infinitely many terms (HW2 P2).

Whether a list is independent can depend on the field. $(1,1,0),(0,1,1),(1,0,1)$ is a basis of $\RR^3$, but in $\FF_2^3$ the three vectors add up to $0$.

\subsection{Dimension}
The \vocab{dimension} of $V$ is the number of vectors in any basis.
Examples: $\dim\FF^n=n$, $\dim\PP_n=n+1$ (basis $1,x,\dots,x^n$), $\dim\FF^{m,n}=mn$.

\begin{proposition*}[LADR 2.38, 2.39, 2.42]
  Suppose $\dim V=n$.
  \begin{itemize}
    \ii Any $n$ independent vectors form a basis. Any $n$ vectors that span form a basis. So you only need to check one of the two conditions.
    \ii A subspace has dimension at most $n$, and it has dimension exactly $n$ only if it is all of $V$.
  \end{itemize}
\end{proposition*}

\begin{theorem*}[LADR 2.43]
  $\dim(U+W)=\dim U+\dim W-\dim(U\cap W)$.

  Two consequences:
  \begin{itemize}
    \ii $U+W$ is direct exactly when $\dim(U+W)=\dim U+\dim W$.
    \ii The intersection cannot be too small: $\dim(U\cap W)\ge\dim U+\dim W-\dim V$.
  \end{itemize}
\end{theorem*}

\traps There is no similar formula for three subspaces.
To show two subspaces are equal, show that one is inside the other and that they have the same dimension. Having the same dimension alone is not enough.

\subsection{Linear maps, null space, range}
\begin{definition*}
  $T\colon V\to W$ is \vocab{linear} if $T(\lambda v+w)=\lambda Tv+Tw$ for all vectors and scalars.

  The \vocab{null space} $\nul T$ is the set of vectors that $T$ sends to $0$.
  The \vocab{range} $\range T$ is the set of all outputs of $T$.
  Both are subspaces.
\end{definition*}

\begin{proposition*}[LADR 3.4]
  A linear map is decided by what it does to a basis.
  If $v_1,\dots,v_n$ is a basis of $V$, you can pick $w_1,\dots,w_n$ in $W$ freely, and there is exactly one linear map with $Tv_1=w_1,\dots,Tv_n=w_n$.
\end{proposition*}
This is how you build a linear map in a proof: pick a convenient basis and say where each basis vector goes.

\begin{proposition*}[LADR 3.15]
  $T$ is injective exactly when its null space is just $\{0\}$.
\end{proposition*}
\begin{proof}
  Suppose $\nul T=\{0\}$ and $Tv=Tw$. Then $T(v-w)=Tv-Tw=0$, so $v-w$ is in the null space. So $v-w=0$, meaning $v=w$.
\end{proof}

\begin{theorem*}[Rank--nullity, LADR 3.21]
  If $V$ is finite-dimensional, then $\dim V=\dim\nul T+\dim\range T$.
\end{theorem*}
\begin{proof}
  Start with a basis $u_1,\dots,u_m$ of the null space, and add vectors $v_1,\dots,v_k$ to make a basis of $V$.
  We show that $Tv_1,\dots,Tv_k$ is a basis of the range. Then $\dim V=m+k=\dim\nul T+\dim\range T$.
  \begin{itemize}
    \ii They span the range: $T$ sends each $u_i$ to $0$, so every output $Tv$ is a combination of $Tv_1,\dots,Tv_k$.
    \ii They are independent: if $\sum a_iTv_i=0$, then $T(\sum a_iv_i)=0$, so $\sum a_iv_i$ is in the null space. That makes it a combination of the $u_j$. Since the $u$'s and $v$'s together are independent, all $a_i=0$.
  \end{itemize}
\end{proof}
Many hard proofs use this same basis: a basis of the null space, extended to a basis of $V$.

\begin{corollary*}
  If $V$ is bigger than $W$ ($\dim V>\dim W$), no linear map $V\to W$ is injective.
  If $V$ is smaller than $W$, no linear map $V\to W$ is surjective.
\end{corollary*}

\begin{proposition*}[Preimage formula, HW3 P1]
  For a subspace $U$ of $W$, the set of vectors that $T$ sends into $U$ has dimension
  \[ \dim T^{-1}(U)=\dim\nul T+\dim(U\cap\range T). \]
\end{proposition*}
\begin{proof}
  Apply rank--nullity to $T$, using only inputs from $T^{-1}(U)$.
  The vectors it sends to $0$ are still $\nul T$.
  Its outputs are exactly the vectors of $U$ that $T$ can reach, which is $U\cap\range T$.
\end{proof}

\begin{proposition*}[HW3 P2]
  If $T^2=0$, then every output of $T$ is sent to $0$ by $T$, so $\range T\subseteq\nul T$.

  If $\nul T=\range T$, then rank--nullity gives $\dim V=2\dim\range T$, so $\dim V$ must be even.
\end{proposition*}

\traps Rank--nullity only says the dimensions add up. It does not say $V$ splits as $\nul T\oplus\range T$. For $N(x,y)=(y,0)$, the null space and the range are the same line.

A linear map can send an independent list to a dependent one (unless it is injective).
It sends a basis to a list that spans the range, but that list might not be a basis.

\subsection{Matrices}
\begin{definition*}
  Pick a basis $v_1,\dots,v_n$ of $V$ and a basis $w_1,\dots,w_m$ of $W$.
  The matrix $\MM(T)$ has $m$ rows and $n$ columns.
  Its $j$th column lists the coefficients you need to write $Tv_j$ in terms of $w_1,\dots,w_m$.
\end{definition*}

\begin{proposition*}[LADR 3.43, 3.76]
  Multiplying matrices is the same as composing maps: $\MM(ST)=\MM(S)\MM(T)$.
  Multiplying a matrix by a coordinate vector is the same as applying the map: $\MM(Tv)=\MM(T)\MM(v)$.
\end{proposition*}

\begin{proposition*}
  $Ax$ is $x_1(\text{column }1)+\dots+x_n(\text{column }n)$.
  So the range of $A$ is the span of its columns, and the rank of $A$ is the dimension of that span.

  The solutions of $Ax=b$, if there are any, are one solution $x_0$ plus anything in the null space.
  Over $\FF_2$, if there is a solution, there are exactly $2^{m-\operatorname{rank}A}$ of them, where $m$ is the number of unknowns.
\end{proposition*}

\traps $\MM(T)$ has $\dim W$ rows and $\dim V$ columns, not the other way round.
The matrix changes if you change the bases.

\subsection{Isomorphisms}
\begin{definition*}
  A linear map $T\colon V\to W$ is an \vocab{isomorphism} (or \vocab{invertible}) if there is a linear map $S\colon W\to V$ that undoes it from both sides: $ST=I$ on $V$ and $TS=I$ on $W$.
  When an isomorphism exists, we write $V\cong W$ and say $V$ and $W$ are isomorphic.
\end{definition*}

\begin{proposition*}[LADR 3.60]
  A map has at most one inverse.
  In fact, if $S_1T=I$ and $TS_2=I$, then $S_1=S_2$.
\end{proposition*}
\begin{proof}
  $S_1=S_1(TS_2)=(S_1T)S_2=S_2$.
\end{proof}

\begin{theorem*}[LADR 3.63]
  $T$ is an isomorphism exactly when it is both injective and surjective.
  You never need to check separately that the inverse is linear. That is automatic.
\end{theorem*}
\begin{proof}
  If $T$ has an inverse: when $Tv=0$, apply $T^{-1}$ to get $v=0$, so $T$ is injective. Every $w$ equals $T(T^{-1}w)$, so $T$ is surjective.

  If $T$ is injective and surjective: each $w$ comes from exactly one $v$. Define $T^{-1}w$ to be that $v$.
  To see it is linear: $T$ sends $\lambda T^{-1}w_1+T^{-1}w_2$ to $\lambda w_1+w_2$. Only one vector goes to $\lambda w_1+w_2$, so $T^{-1}(\lambda w_1+w_2)=\lambda T^{-1}w_1+T^{-1}w_2$.
\end{proof}

\begin{theorem*}[LADR 3.70]
  Two finite-dimensional spaces are isomorphic exactly when they have the same dimension.
\end{theorem*}
\begin{proof}
  If $T$ is an isomorphism, its null space is $\{0\}$ and its range is all of $W$. Rank--nullity gives $\dim V=0+\dim W$.
  If the dimensions are equal, pick a basis of each and let $T$ send the first basis to the second.
\end{proof}

\begin{theorem*}[LADR 3.65]
  If $V$ and $W$ have the same finite dimension, then for $T\colon V\to W$ these are all the same condition:
  $T$ is injective; $T$ is surjective; $T$ is an isomorphism.
\end{theorem*}
\begin{proof}
  By rank--nullity, the null space has dimension $0$ exactly when the range has dimension $\dim V$, which equals $\dim W$.
\end{proof}
This is usually the fastest way to prove something is an isomorphism: show that $Tv=0$ forces $v=0$.

It also proves that equations have solutions. To show $Tp=q$ can be solved for every $q$, show $T$ is injective between spaces of the same dimension. The notes use this for polynomial interpolation (Ex.~34) and for $(x^2+5x+7)p''=q$ (Ex.~35).

\begin{proposition*}[HW4 P3]
  Let $V$ be finite-dimensional and $T\colon V\to V$.
  $T$ is an isomorphism exactly when it sends a basis to a basis. It then sends every basis to a basis.
\end{proposition*}
\begin{proof}
  An isomorphism is injective, so it keeps lists independent, and surjective, so the image of a basis spans.
  Conversely, if $Tv_1,\dots,Tv_n$ is a basis, every vector is an output, so $T$ is surjective and therefore an isomorphism by LADR 3.65.
\end{proof}

\begin{proposition*}[LADR 3D]
  If $S$ and $T$ are invertible, then so is $ST$, and $(ST)^{-1}=T^{-1}S^{-1}$ (the order flips).

  If $ST$ is invertible, then $S$ is surjective and $T$ is injective.
\end{proposition*}
\begin{proof}
  $S$ is surjective: any $v$ equals $S\big(T(ST)^{-1}v\big)$, which is an output of $S$.

  $T$ is injective: if $Tv=0$, then $v=(ST)^{-1}S(Tv)=(ST)^{-1}S(0)=0$.
\end{proof}

\begin{corollary*}[Prop.~33]
  If $V$ is finite-dimensional and $S,T\colon V\to V$, then $ST=I$ implies $TS=I$.
  So in this case you only need to check one side.
\end{corollary*}
\begin{proof}
  $ST=I$ makes $T$ injective. By LADR 3.65, $T$ is invertible. Then $S=S(TT^{-1})=(ST)T^{-1}=T^{-1}$.
\end{proof}

\begin{proposition*}[HW4 P1]
  The $2\times2$ matrix $A=\begin{psmallmatrix}a&b\\c&d\end{psmallmatrix}$, over any field, is invertible exactly when $ad-bc\ne0$. Then
  \[ A^{-1}=\frac{1}{ad-bc}\begin{pmatrix}d&-b\\-c&a\end{pmatrix}. \]
\end{proposition*}
\begin{proof}
  Multiplying out gives $A\begin{psmallmatrix}d&-b\\-c&a\end{psmallmatrix}=(ad-bc)I$. If $ad-bc\ne0$, divide by it.

  If $ad-bc=0$: when $a$ or $b$ is nonzero, $A$ sends the nonzero vector $(b,-a)$ to $0$, so it is not injective. When $a=b=0$, every output has first entry $0$, so $A$ is not surjective.
\end{proof}

\begin{proposition*}[HW4 P4]
  A product has rank at most the rank of each factor: $\operatorname{rank}(BC)\le\min(\operatorname{rank}B,\operatorname{rank}C)$.
\end{proposition*}
\begin{proof}
  Every output of $BC$ is an output of $B$, so $\operatorname{rank}(BC)\le\operatorname{rank}B$.
  The outputs of $BC$ are $B$ applied to the outputs of $C$, and a linear map cannot increase the dimension of a subspace, so $\operatorname{rank}(BC)\le\operatorname{rank}C$.
\end{proof}

\begin{proposition*}[Prop.~26]
  Once bases are fixed, sending $T$ to its matrix is an isomorphism from $\LL(V,W)$ to $\FF^{m,n}$.
  So $\dim\LL(V,W)=(\dim V)(\dim W)$.
\end{proposition*}

\begin{proposition*}[HW4 P7]
  Let $U$ be a subspace of a finite-dimensional $V$, and let $S\colon U\to V$ be injective.
  Then there is an invertible $T\colon V\to V$ that agrees with $S$ on $U$.
\end{proposition*}
\begin{proof}
  Take a basis $u_1,\dots,u_m$ of $U$ and add $w_1,\dots,w_k$ to get a basis of $V$.
  Because $S$ is injective, $Su_1,\dots,Su_m$ is independent. Add $y_1,\dots,y_k$ to it to get a basis of $V$.
  Define $T$ by $Tu_i=Su_i$ and $Tw_j=y_j$. It sends a basis to a basis, so it is invertible.
\end{proof}

\traps Do not write $T^{-1}$ until you have shown it exists.

LADR 3.65 and Prop.~33 need finite dimensions. On $\FF^\infty$, shifting left, $L(a_1,a_2,\dots)=(a_2,a_3,\dots)$, and shifting right, $R(a_1,a_2,\dots)=(0,a_1,a_2,\dots)$, satisfy $LR=I$ but $RL\ne I$. Differentiation on all polynomials is surjective but not injective.

A matrix that is not square has no two-sided inverse, but it can have a one-sided one. $\begin{psmallmatrix}1&2\end{psmallmatrix}$ has infinitely many right inverses.

If $ST$ is injective, you only learn that $T$ is injective, not $S$.

\subsection{Products}
\begin{definition*}
  The \vocab{product} $V_1\times\dots\times V_m$ is the set of lists $(v_1,\dots,v_m)$ with each $v_i$ in $V_i$.
  You add lists entry by entry, and multiply a list by a scalar by multiplying each entry.
\end{definition*}
The spaces do not need to be related. An element of $\PP_2(\RR)\times\RR^3$ is a pair: a polynomial and a vector.
$\RR^m\times\RR^n$ is really the same as $\RR^{m+n}$: just write the two lists one after the other.
Each $V_i$ sits inside the product as the lists that are $0$ except in slot $i$. The product is the direct sum of these copies.

\begin{proposition*}[Prop.~36]
  The dimensions add: $\dim(V_1\times\dots\times V_m)=\dim V_1+\dots+\dim V_m$.
\end{proposition*}
\begin{proof}
  For two spaces: if $v_1,\dots,v_n$ is a basis of $V$ and $w_1,\dots,w_k$ is a basis of $W$, then
  $(v_1,0),\dots,(v_n,0),(0,w_1),\dots,(0,w_k)$ is a basis of $V\times W$.
  It spans because $(v,w)=(v,0)+(0,w)$. It is independent because you can check each slot separately.
\end{proof}

\begin{proposition*}[Prop.~37]
  For subspaces $U_1,\dots,U_m$ of $V$, the map that adds up the entries, $(u_1,\dots,u_m)\mapsto u_1+\dots+u_m$, is a linear map from the product onto the sum $U_1+\dots+U_m$.
  It is injective exactly when the sum is direct.
  So $\dim(U_1+\dots+U_m)\le\dim U_1+\dots+\dim U_m$, with equality exactly when the sum is direct.
\end{proposition*}

\traps The dimension of a product is a sum. The dimension of $\LL(V,W)$ is a product.

\subsection{Affine subsets}
\begin{definition*}
  An \vocab{affine subset} is a subspace moved away from the origin: $v+U$ is the set of all $v+u$ with $u$ in $U$.
  We say $v+U$ is \vocab{parallel} to $U$.
\end{definition*}
Examples:
\begin{itemize}
  \ii The line $x+y=1$ in $\RR^2$ is $(1,0)+\spn((-1,1))$: the line through the origin with slope $-1$, moved to pass through $(1,0)$.
  \ii The antiderivatives of $x^2$ are $\tfrac13x^3+C$ for any constant $C$. That is $\tfrac13x^3$ plus the subspace of constants.
  \ii The solutions of $Tx=b$ are either none, or one solution $x_0$ plus the null space: $x_0+\nul T$.
\end{itemize}
A map of the form $x\mapsto Tx+b$ with $T$ linear is called affine. It is linear only when $b=0$, because it sends $0$ to $b$.

\begin{proposition*}[LADR 3.101]
  For a subspace $U$, these three statements say the same thing:
  \begin{enumerate}
    \ii $v-w$ is in $U$;
    \ii $v+U=w+U$;
    \ii $v+U$ and $w+U$ have at least one point in common.
  \end{enumerate}
  So two shifted copies of $U$ are either identical or do not touch.
\end{proposition*}
\begin{proof}
  If $v-w$ is in $U$, then any $v+u$ can be rewritten as $w+\big((v-w)+u\big)$, which is in $w+U$. The same argument works the other way.

  If the two sets share a point $v+u_1=w+u_2$, then $v-w=u_2-u_1$ is in $U$.
\end{proof}

\begin{proposition*}[HW4 P6, Prop.~38]
  $v+U$ is a subspace exactly when $v$ is in $U$, and then $v+U$ is just $U$.
  Each affine subset is parallel to only one subspace.
\end{proposition*}

\traps The set $v+U$ determines $U$, but not $v$. Any point of the set can play the role of $v$.
The solutions of $Tx=b$ with $b\ne0$ do not form a subspace, because $0$ is not a solution.

\subsection{Quotient spaces}
\begin{definition*}
  Let $U$ be a subspace of $V$. The \vocab{quotient} $V/U$ is the set of all the shifted copies $v+U$.
  You add and scale them using any representative:
  \[ (v+U)+(w+U)=(v+w)+U,\qquad \lambda(v+U)=\lambda v+U. \]
  The zero of $V/U$ is $0+U$, which is $U$ itself.
\end{definition*}
The idea: $V/U$ is $V$ where you stop telling apart two vectors that differ by something in $U$.

Example: take $U$ to be the $y$-axis in $\RR^2$. Then $(1,2)+U$ is the vertical line $x=1$, and so is $(1,4)+U$. Every vertical line is determined by its $x$-value, so $\RR^2/U$ is just a copy of $\RR$.

\textbf{Why ``well defined'' comes up.}
The same element of $V/U$ has many names: $v+U$ and $v'+U$ are equal whenever $v-v'$ is in $U$.
A formula that uses $v$ must give the same answer for every name.
For addition: if $v_1-v_2$ and $w_1-w_2$ are in $U$, then so is $(v_1+w_1)-(v_2+w_2)$, so the answer does not depend on the names. Scalar multiplication works the same way.

\begin{proposition*}[LADR 3.104, 3.105]
  The map $\pi(v)=v+U$ from $V$ to $V/U$ is linear and surjective, and it sends exactly the vectors of $U$ to zero.
  Rank--nullity then gives $\dim V/U=\dim V-\dim U$.

  To get a basis: take a basis $u_1,\dots,u_m$ of $U$ and add $v_1,\dots,v_k$ to make a basis of $V$. Then $v_1+U,\dots,v_k+U$ is a basis of $V/U$.
\end{proposition*}

\begin{proposition*}[HW5 P2]
  Let $T\colon V\to W$ be linear, and suppose $T$ sends everything in $U$ to $0$.
  Then $\tilde T(v+U)=Tv$ is a well-defined linear map from $V/U$ to $W$.
\end{proposition*}
\begin{proof}
  If $v+U=v'+U$, then $v-v'$ is in $U$, so $Tv-Tv'=T(v-v')=0$. Both names give the same output.
\end{proof}

\begin{theorem*}[First isomorphism theorem, LADR 3.107]
  For any linear $T\colon V\to W$, the map $v+\nul T\mapsto Tv$ is an isomorphism from $V/\nul T$ to $\range T$.
\end{theorem*}
\begin{proof}
  It is well defined by the previous proposition.
  It is surjective because every element of the range is some $Tv$.
  It is injective because if $Tv=0$, then $v$ is in the null space, so $v+\nul T$ is the zero element.
\end{proof}
Example: differentiation on polynomials sends exactly the constants to $0$. So ``polynomials, where you ignore the constant term'' is isomorphic to ``polynomials'', by $p\mapsto p'$. The inverse is the indefinite integral.

The second and third isomorphism theorems are $U/(U\cap W)\cong(U+W)/W$ and, for $W\subseteq U$, $(V/W)/(U/W)\cong V/U$. Both are proved the same way: write down a surjective map and apply the first theorem.

\traps Any map you define on $V/U$ by a formula in $v$ needs a check that it is well defined.
For example, $(x,y)+U\mapsto x$ works when $U$ is the $y$-axis, but not when $U$ is the $x$-axis: then $(0,0)+U=(1,0)+U$, and the formula gives $0$ for one and $1$ for the other.

$v+U$ is the zero of $V/U$ when $v$ is in $U$. It does not require $v=0$.

\section{Problems}
Starred problems are the most likely quiz material. Solutions are in Section~\ref{sec:sol}.

\begin{problem}[$\star$ Images and preimages]
  Let $f\colon X\to Y$ be a function.
  \begin{enumerate}[(a)]
    \ii Prove $f$ is injective if and only if $f^{-1}(f(U))=U$ for every $U\subseteq X$.
    \ii Prove $f$ is surjective if and only if $f(f^{-1}(V))=V$ for every $V\subseteq Y$.
    \ii Prove $f(A\cap B)\subseteq f(A)\cap f(B)$, with equality for all $A,B$ if and only if $f$ is injective.
    \ii Let $T\colon V\to W$ be linear, $U\subseteq V$ and $X\subseteq W$ subspaces. Prove
    \[ T^{-1}(T(U))=U+\nul T, \qquad T(T^{-1}(X))=X\cap\range T. \]
    \ii For $V$ finite-dimensional, prove $\dim T(U)=\dim U-\dim(U\cap\nul T)$.
    Then prove $T^{-1}(T(U))=U$ for every subspace $U$ if and only if $T$ is injective.
  \end{enumerate}
\end{problem}

\begin{problem}[Finite fields and polynomial functions]
  Let $p$ be prime. Here $\PP_n(\Fp)$ means formal polynomials of degree $\le n$ (basis $1,x,\dots,x^n$), and $\SF$ is the space of all functions $\Fp\to\Fp$.
  \begin{enumerate}[(a)]
    \ii Prove $\ZZ/6\ZZ$ is not a field.
    \ii Let $\FF$ be a field with $q$ elements and $a\ne0$.
    Prove $x\mapsto ax$ is a bijection $\FF\to\FF$, and deduce $a^{q-1}=1$.
    Conclude $c^p=c$ for all $c\in\Fp$.
    \ii Find $\dim\SF$.
    Let $E\colon\PP_n(\Fp)\to\SF$ send a polynomial to the function it defines.
    Determine exactly for which $n$ the map $E$ is injective, surjective, an isomorphism.
    For $n=p$, give a nonzero element of $\nul E$.
    \ii Over $\FF_2$, expand $(x+1)^4$.
    Which polynomial of degree $\le 1$ defines the same function?
    Which hypothesis of the interpolation argument (Ex.~34) fails when $|\FF|<n+1$?
  \end{enumerate}
\end{problem}

\begin{problem}[Independence depends on the field]
  Let $v_1,\dots,v_n$ be linearly independent, $n\ge2$.
  \begin{enumerate}[(a)]
    \ii Prove or disprove: $v_1-v_2,\ v_2-v_3,\ \dots,\ v_{n-1}-v_n,\ v_n-v_1$ is linearly independent.
    \ii Let $w_j=v_j+v_{j+1}$ for $j<n$ and $w_n=v_n+v_1$.
    Determine exactly, in terms of $n$ and $\FF$, when $w_1,\dots,w_n$ is independent.
    In that case, write $v_1$ in terms of the $w_j$.
    \ii Suppose $v_1+w,\dots,v_n+w$ is linearly dependent. Prove $w\in\spn(v_1,\dots,v_n)$.
    \ii Prove or disprove: if $\spn(u_1,\dots,u_n)=\spn(v_1,\dots,v_n)$, then $u_1,\dots,u_n$ is independent.
  \end{enumerate}
\end{problem}

\begin{problem}[Dimension counting]
  \leavevmode
  \begin{enumerate}[(a)]
    \ii Let $\dim V=10$ and $U,W,X$ subspaces of dimension $7$.
    Prove $U\cap W\cap X\ne\{0\}$, and show $\dim(U\cap W\cap X)=1$ is possible.
    \ii If $U_1,\dots,U_k$ are subspaces of an $n$-dimensional $V$, prove $\dim\bigcap U_i\ge\sum\dim U_i-(k-1)n$.
    \ii Prove or disprove: $\dim(U_1+U_2+U_3)=\sum\dim U_i-\sum_{i<j}\dim(U_i\cap U_j)+\dim(U_1\cap U_2\cap U_3)$.
    \ii Let $\dim V=n$ and $\dim U=m$.
    Prove $\mathcal E=\{T\in\LL(V):T(U)\subseteq U\}$ is a subspace of $\LL(V)$ and compute $\dim\mathcal E$.
  \end{enumerate}
\end{problem}

\begin{problem}[$\star$ Null spaces of compositions]
  Let $T\colon U\to V$ and $S\colon V\to W$ be linear with $U,V$ finite-dimensional.
  \begin{enumerate}[(a)]
    \ii Show $\nul(ST)=T^{-1}(\nul S)$, and deduce
    $\dim\nul(ST)=\dim\nul T+\dim(\nul S\cap\range T)$.
    \ii Prove $\dim\nul(ST)\le\dim\nul S+\dim\nul T$, with equality if and only if $\nul S\subseteq\range T$.
    \ii Prove $\rank(ST)\ge\rank S+\rank T-\dim V$.
    Determine exactly when $\rank(ST)=\rank T$.
    \ii Are there $S,T\in\LL(\RR^3)$ with $\rank S=\rank T=2$ and $ST=0$?
    Are there $S,T\in\LL(\FF^2)$ with $ST\ne0$ and $TS=0$?
  \end{enumerate}
\end{problem}

\begin{problem}[$\star$ Prescribing null spaces and ranges]
  \leavevmode
  \begin{enumerate}[(a)]
    \ii If $T\in\LL(V)$, $V$ finite-dimensional, and $T^2=0$, prove $\rank T\le\frac12\dim V$.
    \ii Prove some $T\in\LL(\FF^n)$ has $\nul T=\range T$ if and only if $n$ is even.
    Write one down for $n=6$.
    \ii Let $V,W$ be finite-dimensional and $U\subseteq V$.
    Prove some $T\in\LL(V,W)$ has $\nul T=U$ if and only if $\dim U\ge\dim V-\dim W$.
    \ii Prove no linear map $\FF^5\to\FF^2$ has null space $\{x: x_1=3x_2,\ x_3=x_4=x_5\}$.
    \ii If $\nul T\cap\range T=\{0\}$ for $T\in\LL(V)$, $V$ finite-dimensional, prove $V=\nul T\oplus\range T$.
  \end{enumerate}
\end{problem}

\begin{problem}[Matrices of polynomial maps]
  \leavevmode
  \begin{enumerate}[(a)]
    \ii Define $\Delta\colon\PP_3(\RR)\to\PP_2(\RR)$ by $(\Delta p)(x)=p(x+1)-p(x)$.
    Explain why $\deg\Delta p\le2$.
    Find $\MM(\Delta)$ in the monomial bases, and find $\nul\Delta$ and $\range\Delta$.
    \ii Find $p$ with $\Delta p=x^2$ and $p(0)=0$. Use it to compute $\sum_{k=0}^{n-1}k^2$.
    \ii Find $\dim\{p\in\PP_3(\RR):\Delta p\in\spn(1,x^2)\}$ using the preimage formula.
    \ii Let $T=I+D$ on $\PP_3(\RR)$, $Dp=p'$.
    Prove $T$ is an isomorphism without a matrix, show $T^{-1}=I-D+D^2-D^3$, and solve $p+p'=x^3$.
    \ii Is $I+D$ invertible on $\PP(\RR)$? Compare with $D$ on $\PP(\RR)$.
  \end{enumerate}
\end{problem}

\begin{problem}[$\star$ Interpolation with derivatives]
  \leavevmode
  \begin{enumerate}[(a)]
    \ii For $a\in\RR$ let $T_a\colon\PP_2(\RR)\to\RR^3$, $T_a(p)=(p(0),p(1),p'(a))$.
    Determine exactly which $a$ make $T_a$ an isomorphism.
    For the others, find $\nul T_a$ and $\range T_a$, and explain using the mean value theorem.
    \ii Prove, before computing, that exactly one $p\in\PP_3(\RR)$ has
    $p(0)=1$, $p'(0)=0$, $p(1)=2$, $p'(1)=-1$. Then find it.
    \ii Prove or disprove: for distinct $x_0,x_1$ and any $a_0,a_1,b_0,b_1$, exactly one $p\in\PP_3(\RR)$ has $p(x_i)=a_i$, $p'(x_i)=b_i$.
  \end{enumerate}
\end{problem}

\begin{problem}[$\star$ Invertibility and one-sided inverses]
  \leavevmode
  \begin{enumerate}[(a)]
    \ii Over $\FF_5$, for which $a$ is $A_a=\begin{psmallmatrix}a&2\\3&a\end{psmallmatrix}$ invertible?
    Find $A_2^{-1}$.
    \ii Over $\CC$, is $B=\begin{psmallmatrix}1+i&2\\1&1-i\end{psmallmatrix}$ invertible?
    If not, find a nonzero vector in $\nul B$ and describe $\range B$.
    \ii Count the invertible $2\times2$ matrices over $\Fp$.
    \ii Let $A=\begin{psmallmatrix}1&2&3\\0&1&1\end{psmallmatrix}\in\RR^{2,3}$.
    Find all $B$ with $AB=I_2$. Prove no $C$ has $CA=I_3$.
    \ii Prove $A\in\FF^{m,n}$ has a right inverse if and only if $\rank A=m$, and a left inverse if and only if $\rank A=n$.
  \end{enumerate}
\end{problem}

\begin{problem}[$\star$ Structure theorems for isomorphisms]
  All spaces are finite-dimensional.
  \begin{enumerate}[(a)]
    \ii Let $T_1,T_2\in\LL(V,W)$.
    Prove $\nul T_1=\nul T_2$ if and only if $T_1=ST_2$ for some invertible $S\in\LL(W)$.
    \ii Prove $\range T_1=\range T_2$ if and only if $T_1=T_2S$ for some invertible $S\in\LL(V)$.
    \ii If $ST=TS$ for every $S\in\LL(V)$, prove $T=\lambda I$.
    \ii Over $\RR$, suppose $T^2-3T+2I=0$.
    Prove $T$ is invertible and write $T^{-1}$ as a polynomial in $T$.
    If $V\ne\{0\}$, prove $T-I$ or $T-2I$ is not invertible.
    What goes wrong over $\FF_2$?
    \ii If $\dim V\ge2$, prove the non-invertible operators on $V$ do not form a subspace.
  \end{enumerate}
\end{problem}

\begin{problem}[$\star$ Products, affine sets, extensions]
  \leavevmode
  \begin{enumerate}[(a)]
    \ii Prove two translates $v+U$, $w+U$ of a subspace $U$ are equal or disjoint.
    \ii Prove $\{x:Tx=b\}$ is empty or equal to $x_0+\nul T$ for any solution $x_0$.
    Use this to list all solutions over $\FF_2$ of $x+y+z+w=1$, $x+z=0$.
    \ii Prove $V\times W=(V\times\{0\})\oplus(\{0\}\times W)$ and $V\cong V\times\{0\}$.
    Reprove $\dim(V\times W)=\dim V+\dim W$.
    \ii Is $\PP_2(\RR)\times\RR^2\cong\RR^{1,5}$? Is $\LL(\RR^2,\RR^3)\cong\RR^2\times\RR^4$?
    \ii Let $U\subseteq V$, $V$ finite-dimensional, and $S\colon U\to V$ injective.
    Prove $S$ extends to an invertible $T\in\LL(V)$.
    Show the hypothesis ``injective'' cannot be dropped.
  \end{enumerate}
\end{problem}

\begin{problem}[Quotient spaces]
  \leavevmode
  \begin{enumerate}[(a)]
    \ii Let $U=\spn((1,1,1))\subseteq\RR^3$.
    Decide whether $(1,0,0)+U=(0,1,-1)+U$ and whether $(2,1,1)+U=(1,0,0)+U$.
    Find a basis of $\RR^3/U$, and write $(3,5,7)+U$ in that basis.
    \ii Prove: if $u_1,\dots,u_m$ is a basis of $U$ and $u_1,\dots,u_m,v_1,\dots,v_k$ is a basis of $V$, then $v_1+U,\dots,v_k+U$ is a basis of $V/U$.
    Conclude $\dim V/U=\dim V-\dim U$ without using rank--nullity.
    \ii On $\RR^2/\spn(e_2)$, show that $f\big((x,y)+U\big)=x$ is a well-defined isomorphism onto $\RR$.
    Show that the same formula is not well defined on $\RR^2/\spn(e_1)$.
    \ii Let $T\in\LL(V,W)$ and $U\subseteq V$ a subspace.
    Prove $\tilde T(v+U)=Tv$ defines a map $V/U\to W$ if and only if $U\subseteq\nul T$.
    In that case, find $\nul\tilde T$.
    \ii Let $T\colon\PP_3(\RR)\to\RR^2$, $T(p)=(p(0),p(1))$.
    Find $\nul T$ and a basis of $\PP_3(\RR)/\nul T$, and describe the isomorphism $\PP_3(\RR)/\nul T\to\RR^2$ from the first isomorphism theorem.
  \end{enumerate}
\end{problem}

\section{True or false}
Give a one-line reason or counterexample. Spaces are finite-dimensional unless stated otherwise.
\begin{enumerate}
  \ii Over $\FF_3$, every nonempty subset of $\FF_3^n$ closed under addition is a subspace.
  \ii Over $\RR$, every nonempty subset of $\RR^n$ closed under addition is a subspace.
  \ii If $U\cup W$ is a subspace, then $U\subseteq W$ or $W\subseteq U$.
  \ii $(1,1,0),(0,1,1),(1,0,1)$ is a basis of $\FF_2^3$.
  \ii There is an injective linear map $\PP_4(\RR)\to\RR^4$.
  \ii If $ST$ is injective, then $S$ is injective.
  \ii If $T^2=0$, then $I-T$ is invertible.
  \ii On an infinite-dimensional space, every injective operator is surjective.
  \ii If $\dim V=5$ and $T\in\LL(V,\FF^3)$, then $\dim\nul T\ge2$.
  \ii For every subspace $U\subseteq V$ there is $T\in\LL(V)$ with $\nul T=\range T=U$.
  \ii If $v_1,\dots,v_n$ is a basis and $T$ is surjective, then $Tv_1,\dots,Tv_n$ is a basis of $W$.
  \ii $\LL(\FF^2,\FF^3)\cong\FF^6$.
  \ii The non-invertible $2\times2$ matrices over $\FF_2$ form a subspace.
  \ii If $S,T\in\LL(V)$ and $ST=I$, then $TS=I$.
  \ii For $A\in\FF^{2,3}$, some $B$ has $AB=I_2$ if and only if the columns of $A$ span $\FF^2$.
\end{enumerate}

\newpage
\section{Solutions}\label{sec:sol}

\subsection*{Problem 1}
\begin{enumerate}[(a)]
  \ii ($\Rightarrow$) If $x\in f^{-1}(f(U))$ then $f(x)=f(u)$ for some $u\in U$, so $x=u\in U$.
  ($\Leftarrow$) If $f(x_1)=f(x_2)$, take $U=\{x_1\}$: then $x_2\in f^{-1}(f(U))=U$, so $x_2=x_1$.
  \ii ($\Rightarrow$) For $y\in V$ pick $x$ with $f(x)=y$; then $x\in f^{-1}(V)$.
  ($\Leftarrow$) Take $V=Y$: $Y=f(f^{-1}(Y))=f(X)$.
  \ii The inclusion is immediate.
  If $f$ is injective and $y=f(a)=f(b)$ with $a\in A$, $b\in B$, then $a=b\in A\cap B$.
  If $f(x_1)=f(x_2)$ with $x_1\ne x_2$, then $A=\{x_1\}$, $B=\{x_2\}$ give $f(A\cap B)=\varnothing\ne f(A)\cap f(B)$.
  \ii $v\in T^{-1}(T(U))\iff Tv=Tu$ for some $u\in U\iff v-u\in\nul T$ for some $u\in U\iff v\in U+\nul T$.
  And $w\in T(T^{-1}(X))\iff w=Tv$ with $Tv\in X\iff w\in\range T\cap X$.
  \ii $T|_U$ has null space $U\cap\nul T$ and range $T(U)$; apply rank--nullity.
  By (d), $T^{-1}(T(U))=U\iff\nul T\subseteq U$.
  This holds for $U=\{0\}$ if and only if $\nul T=\{0\}$.
\end{enumerate}

\subsection*{Problem 2}
\begin{enumerate}[(a)]
  \ii $2\cdot3=0$ with $2,3\ne0$, impossible in a field.
  \ii $ax=ay\implies a(x-y)=0\implies x=y$.
  An injective self-map of a finite set is bijective, and it fixes $0$, so it permutes the nonzero $x_1,\dots,x_{q-1}$.
  Then $a^{q-1}\prod x_i=\prod(ax_i)=\prod x_i$, and $\prod x_i\ne0$ cancels: $a^{q-1}=1$.
  So $a^q=a$ for all $a$ (including $0$).
  \ii The indicator functions $\delta_c$ form a basis, so $\dim\SF=p$; $\dim\PP_n=n+1$.
  \begin{itemize}
    \ii Injective if and only if $n\le p-1$: a nonzero polynomial of degree $\le n<p$ cannot vanish at all $p$ points.
    If $n\ge p$, then $x^p-x\in\nul E$ by (b).
    \ii Surjective if and only if $n\ge p-1$: at $n=p-1$, $E$ is injective between $p$-dimensional spaces, so surjective; larger $n$ only enlarges the range.
    For $n<p-1$ the domain is too small.
    \ii Isomorphism if and only if $n=p-1$.
  \end{itemize}
  \ii $(x+1)^4=(x^2+1)^2=x^4+1$, which equals $x+1$ as a function since $c^4=c$ on $\FF_2$.
  Ex.~34 needs $n+1$ distinct points in $\FF$; with $|\FF|<n+1$ they do not exist, and distinct polynomials (e.g.\ $x^2$ and $x$ over $\FF_2$) define the same function.
\end{enumerate}

\subsection*{Problem 3}
\begin{enumerate}[(a)]
  \ii False: the vectors sum to $0$.
  \ii If $\sum c_jw_j=0$, the coefficient of $v_j$ is $c_j+c_{j-1}$ (indices mod $n$), so $c_j=(-1)^{j-1}c_1$, and wrapping around gives $c_1=(-1)^nc_1$.
  \begin{itemize}
    \ii $n$ even: $c=(1,-1,\dots,-1)$ works, so the list is dependent over every field.
    \ii $n$ odd: $2c_1=0$. Independent if and only if $2\ne0$ in $\FF$; dependent over $\FF_2$.
  \end{itemize}
  When independent, $w_1-w_2+w_3-\dots+w_n=2v_1$, so $v_1=\tfrac12(w_1-w_2+\dots+w_n)$.
  \ii Take $\sum a_i(v_i+w)=0$, not all $a_i=0$.
  Then $\sum a_iv_i=-(\sum a_i)w$.
  If $s=\sum a_i=0$, independence forces all $a_i=0$, contradiction.
  So $w=-s^{-1}\sum a_iv_i$.
  \ii True: $\spn(v_1,\dots,v_n)$ has dimension $n$, and a spanning list of length $n$ in it is a basis.
\end{enumerate}

\subsection*{Problem 4}
\begin{enumerate}[(a)]
  \ii $\dim(U\cap W)\ge7+7-10=4$, then $\dim(U\cap W\cap X)\ge4+7-10=1$.
  Equality: in $\FF^{10}$, $U=\{x_1=x_2=x_3=0\}$, $W=\{x_4=x_5=x_6=0\}$, $X=\{x_7=x_8=x_9=0\}$ meet in $\spn(e_{10})$.
  \ii Induct with $\dim(I\cap U_k)\ge\dim I+\dim U_k-n$ where $I=U_1\cap\dots\cap U_{k-1}$.
  \ii False: lines spanned by $e_1,e_2,e_1+e_2$ in $\RR^2$ give $2$ on the left and $3$ on the right.
  \ii Closure: $(\lambda S+T)u=\lambda Su+Tu\in U$.
  Extend a basis $u_1,\dots,u_m$ of $U$ to a basis of $V$.
  $T(U)\subseteq U$ if and only if rows $m+1,\dots,n$ of columns $1,\dots,m$ of $\MM(T)$ vanish.
  Since $\MM$ is an isomorphism, $\dim\mathcal E=n^2-m(n-m)$.
\end{enumerate}

\subsection*{Problem 5}
\begin{enumerate}[(a)]
  \ii $u\in\nul(ST)\iff Tu\in\nul S$. Apply HW3 P1(b) with $U=\nul S$.
  \ii $\nul S\cap\range T\subseteq\nul S$ gives the inequality.
  Equality if and only if the two nested subspaces are equal, if and only if $\nul S\subseteq\range T$.
  \ii $\rank(ST)=\dim U-\dim\nul(ST)\ge\dim U-\dim\nul T-\dim\nul S=\rank T+\rank S-\dim V$.
  Exactly: $\rank(ST)=\rank T-\dim(\nul S\cap\range T)$, so $\rank(ST)=\rank T$ if and only if $\nul S\cap\range T=\{0\}$.
  \ii No: $\rank(ST)\ge2+2-3=1$.
  Yes: $S=\begin{psmallmatrix}1&0\\0&0\end{psmallmatrix}$, $T=\begin{psmallmatrix}0&1\\0&0\end{psmallmatrix}$ give $ST=T\ne0$ and $TS=0$.
\end{enumerate}

\subsection*{Problem 6}
\begin{enumerate}[(a)]
  \ii $Tw=T^2v=0$ for $w=Tv$, so $\range T\subseteq\nul T$, and $\dim V\ge2\rank T$.
  \ii ($\Rightarrow$) $n=2\rank T$.
  ($\Leftarrow$) For $n=2k$ set $Te_j=0$ and $Te_{k+j}=e_j$ ($j\le k$).
  Then $\range T=\spn(e_1,\dots,e_k)\subseteq\nul T$ and both have dimension $k$.
  For $n=6$: $\MM(T)=\begin{psmallmatrix}0&I_3\\0&0\end{psmallmatrix}$.
  \ii ($\Rightarrow$) $\dim V-\dim U=\rank T\le\dim W$.
  ($\Leftarrow$) Extend a basis $u_1,\dots,u_m$ of $U$ by $v_1,\dots,v_k$, where $k\le\dim W$.
  Set $Tu_i=0$ and $Tv_i=w_i$ for independent $w_1,\dots,w_k$.
  If $T(\sum a_iu_i+\sum b_iv_i)=\sum b_iw_i=0$, all $b_i=0$, so $\nul T=U$.
  \ii The subspace $\{(3s,s,t,t,t)\}$ has dimension $2<5-2$.
  \ii The sum is direct and has dimension $\dim\nul T+\dim\range T=\dim V$.
\end{enumerate}

\subsection*{Problem 7}
\begin{enumerate}[(a)]
  \ii The $x^k$ terms of $(x+1)^k-x^k$ cancel.
  $\Delta1=0$, $\Delta x=1$, $\Delta x^2=2x+1$, $\Delta x^3=3x^2+3x+1$, so
  \[ \MM(\Delta)=\begin{pmatrix}0&1&1&1\\0&0&2&3\\0&0&0&3\end{pmatrix}. \]
  Rank $3$, so $\range\Delta=\PP_2(\RR)$ and $\nul\Delta=\spn(1)$.
  \ii Solving gives $c_3=\tfrac13$, $c_2=-\tfrac12$, $c_1=\tfrac16$:
  $p=\tfrac13x^3-\tfrac12x^2+\tfrac16x$.
  Telescoping, $\sum_{k=0}^{n-1}k^2=p(n)-p(0)=\tfrac{n(n-1)(2n-1)}{6}$.
  \ii $1+\dim(\spn(1,x^2)\cap\PP_2)=1+2=3$.
  \ii If $p\ne0$ has degree $d$, then $p+p'$ has degree $d$, so $T$ is injective, hence an isomorphism.
  $D^4=0$ on $\PP_3$, so $(I+D)(I-D+D^2-D^3)=I+D^4=I$.
  Then $p=(I-D+D^2-D^3)x^3=x^3-3x^2+6x-6$.
  \ii Yes: $Sp=\sum_{k\ge0}(-D)^kp$ is a finite sum for each $p$, and $(I+D)S=S(I+D)=I$.
  $D$ itself is surjective but not injective on $\PP(\RR)$.
\end{enumerate}

\subsection*{Problem 8}
\begin{enumerate}[(a)]
  \ii For $p=c_0+c_1x+c_2x^2$, $T_ap=0$ means $c_0=0$, $c_1+c_2=0$, $c_1+2ac_2=0$, so $(2a-1)c_2=0$.
  For $a\ne\tfrac12$, $T_a$ is injective, hence an isomorphism.
  For $a=\tfrac12$: $\nul T_a=\spn(x-x^2)$, and since $p(1)-p(0)=c_1+c_2=p'(\tfrac12)$,
  $\range T_a=\{(y_0,y_1,z):z=y_1-y_0\}$.
  For a quadratic, the secant slope on $[0,1]$ equals the derivative at the midpoint exactly, so $p'(\tfrac12)$ cannot be prescribed.
  \ii $Sp=(p(0),p'(0),p(1),p'(1))$ on $\PP_3\to\RR^4$.
  If $Sp=0$ then $p=c_2x^2+c_3x^3$ with $c_2+c_3=0=2c_2+3c_3$, so $p=0$.
  Hence $S$ is an isomorphism.
  Solving: $p=1+4x^2-3x^3$.
  \ii True: a $p\in\PP_3$ with double roots at $x_0,x_1$ is divisible by the degree-$4$ polynomial $(x-x_0)^2(x-x_1)^2$, so $p=0$.
\end{enumerate}

\subsection*{Problem 9}
\begin{enumerate}[(a)]
  \ii $\det A_a=a^2-6=(a-1)(a+1)$ in $\FF_5$, so invertible if and only if $a\in\{0,2,3\}$.
  For $a=2$, $\det=3$ and $3^{-1}=2$:
  $A_2^{-1}=2\begin{psmallmatrix}2&-2\\-3&2\end{psmallmatrix}=\begin{psmallmatrix}4&1\\4&4\end{psmallmatrix}$.
  \ii $\det B=(1+i)(1-i)-2=0$.
  $(2,-1-i)^T\in\nul B$.
  Column 2 is $(1-i)$ times column 1, so $\range B=\spn\big((1+i,1)^T\big)$.
  \ii $(p^2-1)(p^2-p)$: any nonzero first column, then a second column that is not a multiple of it.
  \ii $\nul A=\spn\big((-1,-1,1)^T\big)$; particular solutions $(1,0,0)^T$ for $e_1$ and $(-2,1,0)^T$ for $e_2$. So
  \[ B=\begin{pmatrix}1-s&-2-t\\-s&1-t\\s&t\end{pmatrix},\quad s,t\in\RR. \]
  No left inverse: $\rank(CA)\le\rank A\le2<3$.
  \ii $AB=I_m\implies m\le\rank A\le m$.
  Conversely, if $A$ is surjective, choose $b_j$ with $Ab_j=e_j$.
  $CA=I_n\implies n\le\rank A\le n$.
  Conversely, if $A$ is injective, extend $Ae_1,\dots,Ae_n$ to a basis of $\FF^m$ and let $C$ send $Ae_j\mapsto e_j$ and the rest to $0$.
\end{enumerate}

\subsection*{Problem 10}
\begin{enumerate}[(a)]
  \ii ($\Leftarrow$) $S$ is injective, so $ST_2v=0\iff T_2v=0$.
  ($\Rightarrow$) Take a basis $u_1,\dots,u_m$ of the common null space, extended by $v_1,\dots,v_k$.
  Both $T_1v_1,\dots,T_1v_k$ and $T_2v_1,\dots,T_2v_k$ are independent.
  Extend them to bases $T_2v_i,x_j$ and $T_1v_i,y_j$ of $W$ (same length).
  Define $S(T_2v_i)=T_1v_i$ and $Sx_j=y_j$: basis to basis, so $S$ is invertible, and $ST_2=T_1$ on $u_i$ and $v_i$.
  \ii ($\Leftarrow$) $\range(T_2S)=T_2(V)$ since $S$ is surjective.
  ($\Rightarrow$) Equal ranks give equal nullities $m$.
  Take an adapted basis $u_i,v_i$ for $T_1$, choose $v_i'$ with $T_2v_i'=T_1v_i$, and a basis $u_i'$ of $\nul T_2$.
  Then $u_i',v_i'$ is independent (apply $T_2$ to a relation) and has length $\dim V$.
  Set $Su_i=u_i'$ and $Sv_i=v_i'$.
  \ii If $v,Tv$ were independent, extend to a basis and let $Sv=v$, $S(Tv)=0$: then $STv=0\ne Tv=TSv$.
  So $Tv=a_vv$ for every $v$.
  For independent $v,w$: $a_{v+w}(v+w)=a_vv+a_ww$ gives $a_v=a_w$.
  For $w=cv$: $a_ww=T(cv)=a_vw$.
  So $T=\lambda I$.
  \ii $T(3I-T)=2I$, so $T^{-1}=\tfrac12(3I-T)$.
  $(T-I)(T-2I)=0$ is not invertible when $V\ne\{0\}$, so the factors cannot both be.
  Over $\FF_2$ the equation is $T^2-T=0$; dividing by $2$ is impossible and $T=0$ satisfies it.
  \ii With a basis $v_1,\dots,v_n$, let $P_1v_1=v_1$, $P_1v_j=0$ for $j\ge2$, and $P_2=I-P_1$.
  Neither is injective, but $P_1+P_2=I$.
\end{enumerate}

\subsection*{Problem 11}
\begin{enumerate}[(a)]
  \ii If $v+u_1=w+u_2$, then $v-w\in U$, so $v+U=w+U$ (HW4 P6b).
  \ii $Tx=b\iff T(x-x_0)=0$.
  Over $\FF_2$: $z=x$, then $y+w=1$. With $x_0=(0,0,0,1)$ and $\nul=\spn((1,0,1,0),(0,1,0,1))$, the solutions are
  $(0,0,0,1)$, $(1,0,1,1)$, $(0,1,0,0)$, $(1,1,1,0)$.
  \ii $(v,w)=(v,0)+(0,w)$ and the intersection is $\{0\}$.
  $v\mapsto(v,0)$ is a linear bijection onto $V\times\{0\}$.
  Dimensions of a direct sum of two subspaces add.
  \ii Both yes: dimensions $5=5$ and $6=6$.
  \ii Extend a basis $u_i$ of $U$ by $w_j$ to a basis of $V$.
  Injectivity makes $Su_i$ independent; extend by $y_j$ to a basis of $V$.
  Set $Tu_i=Su_i$, $Tw_j=y_j$.
  Counterexample without injectivity: $S=0$ on $U=V=\FF$.
\end{enumerate}

\subsection*{Problem 12}
\begin{enumerate}[(a)]
  \ii $(1,0,0)-(0,1,-1)=(1,-1,1)\notin U$, so the translates are different (and disjoint).
  $(2,1,1)-(1,0,0)=(1,1,1)\in U$, so they are equal.
  $(1,1,1),e_1,e_2$ is a basis of $\RR^3$, so $e_1+U,e_2+U$ is a basis of $\RR^3/U$ by (b).
  $(3,5,7)=7(1,1,1)-4e_1-2e_2$, so $(3,5,7)+U=-4(e_1+U)-2(e_2+U)$.
  \ii Spanning: write $v=\sum a_iu_i+\sum b_jv_j$; then $v+U=\sum b_j(v_j+U)$ because $\sum a_iu_i\in U$.
  Independence: if $\sum b_j(v_j+U)=0+U$, then $\sum b_jv_j\in U$, so $\sum b_jv_j=\sum a_iu_i$ for some $a_i$.
  Independence of the full basis forces all $b_j=0$.
  So $\dim V/U=k=\dim V-\dim U$.
  \ii If $(x,y)+U=(x',y')+U$, then $(x-x',y-y')\in\spn(e_2)$, so $x=x'$: well defined.
  It is linear, surjective ($x\mapsto(x,0)+U$), and injective ($x=0$ means $(0,y)\in U$).
  On $\RR^2/\spn(e_1)$: $(0,0)+U=(1,0)+U$ but the formula gives $0$ and $1$.
  \ii ($\Leftarrow$) If $v-v'\in U\subseteq\nul T$, then $Tv=Tv'$.
  Linearity follows from linearity of $T$.
  ($\Rightarrow$) If $u\in U$ has $Tu\ne0$, then $0+U=u+U$ would have to map to both $0$ and $Tu$.
  $\nul\tilde T=\{v+U:v\in\nul T\}$, which is $(\nul T)/U$.
  In particular $\tilde T$ is injective if and only if $U=\nul T$.
  \ii $T$ is surjective ($p=a+(b-a)x$ hits $(a,b)$), so $\dim\nul T=4-2=2$ and $\nul T=\spn\big(x(x-1),\,x^2(x-1)\big)$.
  $1,x,x(x-1),x^2(x-1)$ is a basis of $\PP_3(\RR)$, so $1+\nul T,\ x+\nul T$ is a basis of the quotient.
  The isomorphism is $p+\nul T\mapsto(p(0),p(1))$, sending $1+\nul T\mapsto(1,1)$ and $x+\nul T\mapsto(0,1)$.
  Two polynomials are equal in the quotient if and only if they agree at $0$ and $1$.
\end{enumerate}

\subsection*{True or false}
\begin{enumerate}
  \ii True: $2v=v+v$ and $0=3v$ lie in the set, and those are all the scalar multiples.
  \ii False: $\ZZ^n$.
  \ii True: if $u\in U\setminus W$, $w\in W\setminus U$, then $u+w$ lies in neither.
  \ii False: the sum is $0$. (True over $\RR$.)
  \ii False: $5>4$.
  \ii False: $T(t)=(t,0)$, $S(x,y)=x$ give $ST=I$. ($T$ must be injective.)
  \ii True: $(I-T)^{-1}=I+T$.
  \ii False: the right shift.
  \ii True: $5-\rank T\ge5-3$.
  \ii False: requires $\dim U=\frac12\dim V$.
  \ii False: spans but need not be independent; $(x,y)\mapsto x$.
  \ii True: $\dim=6$.
  \ii False: $E_{11}+E_{22}=I$.
  \ii True in finite dimensions; false for the shifts on $\FF^\infty$.
  \ii True: a right inverse exists if and only if $\rank A=2$.
\end{enumerate}

\end{document}
